% !TeX program = xelatex
% ============================================================================
%  第 5 课　三角形的对称与置换记号　—— 学生版讲义
%  与 02-课件/第05课-三角形的对称与置换记号/第05课-三角形的对称与置换记号.tex 同步
%  约定：顶点 A 在上、B 在右下、C 在左下；转 120° 指顺时针
% ============================================================================

\zhipianye{5}

\xhlesson{5}{三角形的对称与置换记号}{}

\begin{mubiao}
  \begin{itemize}
    \item 能说出正三角形的 \(6\) 个对称动作。
    \item 会把每个动作写成一张两行式（置换记号）。
    \item 知道两行式的作用：不画图也能说清谁去了哪。
  \end{itemize}
\end{mubiao}

\section*{一、回顾}

\begin{sikao}
  上节课的两个动作：甲 \(=\) 乘 \(2\)，乙 \(=\) 加 \(1\)。取 \(x=1\)，先甲后乙得 \underline{\hspace{2em}}；
  先乙后甲得 \underline{\hspace{2em}}。一样吗？\underline{\hspace{2em}}
\end{sikao}

\section*{二、三角形的六个动作}

\noindent 老师手里有一张正三角形纸片：正面三个顶点写着 \(A,B,C\)，背面是另一种颜色。
约定位置：\(A\) 在上，\(B\) 在右下，\(C\) 在左下。

\begin{definition}[对称动作]
  做完一次动作，图形要和原来的位置\textbf{完全重合}——顶点对准顶点，边对准边。
\end{definition}

% 待插入 \imgfig[0.5\textwidth]{zhipian_sanjiao}：一张正三角形纸片示意（正面 A 在上、B 在右下、C 在左下，背面另一种颜色）
\begin{dongshou}[边看边记]
  老师每做一个动作，就把三角形贴回原位。\\
  这张表\textbf{只记 \(A\) 和 \(C\) 去哪}——第三个顶点 \(B\) 的位置，由 \(A\)、\(C\) 两个位置就能推出来。
  前三个动作跟着老师做、自己填；后三个翻折的动作先空着，等老师翻开这一页\textbf{对答案}。

  \begin{center}
  \begin{tabular}{@{}lc@{}}
    \toprule
    动作 & 顶点去了哪（只记 \(A\) 和 \(C\)） \\
    \midrule
    不动 & \(A\to\underline{\hspace{1.2em}}\)，\(C\to\underline{\hspace{1.2em}}\) \\[4pt]
    顺时针转 \(120^\circ\) & \(A\to\underline{\hspace{1.2em}}\)，\(C\to\underline{\hspace{1.2em}}\) \\[4pt]
    顺时针转 \(240^\circ\) & \(A\to\underline{\hspace{1.2em}}\)，\(C\to\underline{\hspace{1.2em}}\) \\[4pt]
    \midrule
    沿过 \(A\) 的轴翻（对答案） & \(A\to A\)，\(C\to B\) \\[4pt]
    沿过 \(B\) 的轴翻（对答案） & \(A\to C\)，\(C\to A\) \\[4pt]
    沿过 \(C\) 的轴翻（对答案） & \(A\to B\)，\(C\to C\) \\
    \bottomrule
  \end{tabular}
  \end{center}
  （完整写法：沿过 \(A\) 的轴翻是 \(A\to A,\ B\to C,\ C\to B\)；沿过 \(B\) 的轴翻是 \(A\to C,\ B\to B,\ C\to A\)。）
  自查：写完两行式的下排，看\textbf{下排是不是还让 \(A,B,C\) 各出现一次}。
\end{dongshou}

\begin{sikao}
  一共数出了几个动作？\underline{\hspace{2em}} 个。其中转动几个、翻折几个？
  \liukong{2}
\end{sikao}

\section*{三、置换记号}

\begin{definition}[置换的双行记号]
  上排写原来的位置，下排写这个位置上的东西\textbf{去了哪}。两张表合在一起，
  写成一对括号里的两行，就记下了整个动作。
\end{definition}

\begin{example}
  「顺时针转 \(120^\circ\)」写成
  \[
    \begin{pmatrix}
      A & B & C\\
      B & C & A
    \end{pmatrix},
  \]
  读作「\(A\) 去 \(B\)，\(B\) 去 \(C\)，\(C\) 去 \(A\)」。
\end{example}

\begin{example}[翻折也照这样来一遍：沿过 \(A\) 的轴翻]
  过 \(A\) 的那条轴，就是从上顶点 \(A\) 到对边中点的那条线。翻的时候：
  \begin{itemize}
    \item \(A\) 正好在这条轴上，所以 \(A\) 不动，\(A\to A\)；
    \item \(B\) 和 \(C\) 在这条轴的两边，翻过去就\textbf{互相对调}：\(B\to C\)，\(C\to B\)。
  \end{itemize}
  把三个答案排成两行式：
  \[
    \begin{pmatrix}
      A & B & C\\
      A & C & B
    \end{pmatrix},
  \]
  读作「\(A\) 不动，\(B\) 去 \(C\)，\(C\) 去 \(B\)」。
\end{example}

\begin{dongshou}[跟做，写在下面]
  \begin{enumerate}
    \item 顺时针转 \(240^\circ\)：
      \[
        \begin{pmatrix}
          A & B & C\\[2pt]
          \hspace{1.8em} & \hspace{1.8em} & \hspace{1.8em}
        \end{pmatrix}
      \]
    \item 沿过 \(B\) 的轴翻折：
      \[
        \begin{pmatrix}
          A & B & C\\[2pt]
          \hspace{1.8em} & \hspace{1.8em} & \hspace{1.8em}
        \end{pmatrix}
      \]
  \end{enumerate}
  提示：先问「\(A\) 去哪了」，再问「\(B\) 去哪了」，最后问「\(C\) 去哪了」。
  写完自查一遍：\textbf{下排是不是还让 \(A,B,C\) 各出现一次？}
\end{dongshou}

\begin{example}[对一下答案]
  \begin{center}
  \begin{tabular}{@{}lc@{}}
    \toprule
    动作 & 两行式 \\
    \midrule
    不动 & \(\begin{pmatrix}A&B&C\\A&B&C\end{pmatrix}\) \\[8pt]
    顺时针转 \(120^\circ\) & \(\begin{pmatrix}A&B&C\\B&C&A\end{pmatrix}\) \\[8pt]
    顺时针转 \(240^\circ\) & \(\begin{pmatrix}A&B&C\\C&A&B\end{pmatrix}\) \\[8pt]
    沿过 \(A\) 的轴翻 & \(\begin{pmatrix}A&B&C\\A&C&B\end{pmatrix}\) \\[8pt]
    沿过 \(B\) 的轴翻 & \(\begin{pmatrix}A&B&C\\C&B&A\end{pmatrix}\) \\[8pt]
    沿过 \(C\) 的轴翻 & \(\begin{pmatrix}A&B&C\\B&A&C\end{pmatrix}\) \\
    \bottomrule
  \end{tabular}
  \end{center}
\end{example}

\begin{sikao}
  六张两行式里，哪一张的下排和上排一模一样？它表示哪个动作？
  \liukong{2}
\end{sikao}

\begin{sikao}
  反过来考一考：\(\begin{pmatrix}A&B&C\\B&A&C\end{pmatrix}\) 表示哪个动作？\\
  做这个动作的时候，纸片要不要翻面？
  \liukong{2}
\end{sikao}

\begin{xiaojie}
  用自己的话说说：两行式是用来干什么的？
  \liukong{4}
\end{xiaojie}

\noindent\textbf{下节课}：把「先做这个，再做那个」当成一种运算，看看这 \(6\) 个动作能不能凑成一个群。